% Options for packages loaded elsewhere
\PassOptionsToPackage{unicode}{hyperref}
\PassOptionsToPackage{hyphens}{url}
%
\documentclass[
]{article}
\usepackage{amsmath,amssymb}
\usepackage{iftex}
\ifPDFTeX
  \usepackage[T1]{fontenc}
  \usepackage[utf8]{inputenc}
  \usepackage{textcomp} % provide euro and other symbols
\else % if luatex or xetex
  \usepackage{unicode-math} % this also loads fontspec
  \defaultfontfeatures{Scale=MatchLowercase}
  \defaultfontfeatures[\rmfamily]{Ligatures=TeX,Scale=1}
\fi
\usepackage{lmodern}
\ifPDFTeX\else
  % xetex/luatex font selection
\fi
% Use upquote if available, for straight quotes in verbatim environments
\IfFileExists{upquote.sty}{\usepackage{upquote}}{}
\IfFileExists{microtype.sty}{% use microtype if available
  \usepackage[]{microtype}
  \UseMicrotypeSet[protrusion]{basicmath} % disable protrusion for tt fonts
}{}
\makeatletter
\@ifundefined{KOMAClassName}{% if non-KOMA class
  \IfFileExists{parskip.sty}{%
    \usepackage{parskip}
  }{% else
    \setlength{\parindent}{0pt}
    \setlength{\parskip}{6pt plus 2pt minus 1pt}}
}{% if KOMA class
  \KOMAoptions{parskip=half}}
\makeatother
\usepackage{xcolor}
\usepackage{longtable,booktabs,array}
\usepackage{calc} % for calculating minipage widths
% Correct order of tables after \paragraph or \subparagraph
\usepackage{etoolbox}
\makeatletter
\patchcmd\longtable{\par}{\if@noskipsec\mbox{}\fi\par}{}{}
\makeatother
% Allow footnotes in longtable head/foot
\IfFileExists{footnotehyper.sty}{\usepackage{footnotehyper}}{\usepackage{footnote}}
\makesavenoteenv{longtable}
\setlength{\emergencystretch}{3em} % prevent overfull lines
\providecommand{\tightlist}{%
  \setlength{\itemsep}{0pt}\setlength{\parskip}{0pt}}
\setcounter{secnumdepth}{-\maxdimen} % remove section numbering
\ifLuaTeX
  \usepackage{selnolig}  % disable illegal ligatures
\fi
\IfFileExists{bookmark.sty}{\usepackage{bookmark}}{\usepackage{hyperref}}
\IfFileExists{xurl.sty}{\usepackage{xurl}}{} % add URL line breaks if available
\urlstyle{same}
\hypersetup{
  hidelinks,
  pdfcreator={LaTeX via pandoc}}

\author{}
\date{}

\begin{document}

\hypertarget{data-driven-detection-of-complex-multiplication-in-weight-2-cusp-forms}{%
\section{Data-Driven Detection of Complex Multiplication in Weight 2
Cusp
Forms}\label{data-driven-detection-of-complex-multiplication-in-weight-2-cusp-forms}}

\begin{quote}
\textbf{Authors}: Tobias Faller

\textbf{Date}: June 2026

\textbf{Status}: arXiv submission
\end{quote}

\hypertarget{abstract}{%
\subsection{Abstract}\label{abstract}}

We introduce a machine learning approach for detecting Complex
Multiplication (CM) in weight 2 newforms using a dataset of 53,779
modular forms from the LMFDB. By combining prime-indexed Fourier
coefficients a\_p for 25 primes up to 97 with 11 Sato-Tate moments
M\_2(d) and standardized ratios, we achieve F1=0.900 and precision=0.973
on an 80/20 held-out test set using Gradient Boosting Machines (GBM).
Our contribution reveals $M_4$/$M_2$ as the most discriminative feature
(importance 0.157), with trace coefficients at p=23, 41, and 7
contributing significantly. We find CM forms represent only 0.40\% of
the dataset (213/53,779), presenting a challenging class imbalance
problem. Our results demonstrate that CM is learnable from
small-dimensional feature sets without feature selection, providing a
scalable alternative to Elliptic Curve analysis.

\hypertarget{introduction}{%
\subsection{1. Introduction}\label{introduction}}

\hypertarget{background}{%
\subsubsection{1.1 Background}\label{background}}

Complex Multiplication (CM) is a fundamental property of elliptic curves
with deep connections to number theory. An elliptic curve E over $\mathbb{Q}$ has
CM if its endomorphism ring End(E) strictly contains $\mathbb{Z}$, which occurs
precisely when E has complex multiplication by the ring of integers of
an imaginary quadratic field. Classically, CM detection requires:

\begin{enumerate}
\def\labelenumi{\arabic{enumi}.}
\tightlist
\item
  Computing the conductor N of an elliptic curve E/$\mathbb{Q}$
\item
  Testing the j-invariant j(E) against Hilbert class polynomials
\item
  Checking the property that all L-series coefficients a\_p are
  algebraic integers in the CM field
\end{enumerate}

This approach is computationally expensive and requires intimate
knowledge of Elliptic Curve theory. Our work explores whether CM can be
detected directly from the Fourier coefficients of modular forms,
bypassing Elliptic Curve analysis entirely.

\hypertarget{motivation}{%
\subsubsection{1.2 Motivation}\label{motivation}}

The Langlands program connects modular forms to algebraic geometry. In
particular, weight 2 newforms correspond to elliptic curves over $\mathbb{Q}$ via
the modularity theorem. Thus, detecting CM in modular forms is
equivalent to detecting CM in elliptic curves.

Our motivation is threefold:

\begin{enumerate}
\def\labelenumi{\arabic{enumi}.}
\item
  \textbf{Scalability}: The LMFDB now contains tens of thousands of
  modular forms. We seek methods that scale to this data volume without
  per-form specialized analysis.
\item
  \textbf{Classification Interpretability}: Machine learning models
  trained on this data can reveal which Fourier coefficients carry CM
  information, improving theoretical understanding.
\item
  \textbf{Generalization to Higher Weight}: CM detection techniques
  based on Fourier coefficients can potentially generalize to weight
  \textgreater2 newforms, where the Elliptic Curve correspondence fails.
\end{enumerate}

\hypertarget{contributions}{%
\subsubsection{1.3 Contributions}\label{contributions}}

\begin{itemize}
\tightlist
\item
  We build the first large-scale ML dataset for CM detection with 53,779
  labeled weight 2 newforms
\item
  We achieve F1=0.900 on held-out test data using a 36-dimensional
  feature set (25 traces + 11 Sato-Tate moments)
\item
  We identify $M_4$/$M_2$ (importance 0.157) and specific trace coefficients
  (p=23, 41, 7) as the top predictive features
\item
  We validate the class imbalance problem: CM forms constitute only
  0.40\% of the dataset (213/53,779)
\item
  We demonstrate robust performance with only 8 misclassifications out
  of 10,756 test examples (0.074\% error rate)
\end{itemize}

\hypertarget{data-and-methods}{%
\subsection{2. Data and Methods}\label{data-and-methods}}

\hypertarget{dataset-acquisition}{%
\subsubsection{2.1 Dataset Acquisition}\label{dataset-acquisition}}

We extracted weight 2 newforms from the LMFDB PostgreSQL mirror at
\texttt{devmirror.lmfdb.xyz:5432} using the SQL schema documented in
\texttt{scripts/collect\_lmfdb\_sql.py}. The query filters:

\begin{itemize}
\tightlist
\item
  Weight: 2
\item
  Labelarity: 0 (non-cuspidal excluded)
\item
  Character: trivial
\item
  CM labels extracted, computed \texttt{aps\ =\ denotes\_cm(label)} and
  stored manually
\end{itemize}

The resulting dataset contains:

\begin{itemize}
\tightlist
\item
  Total forms: 53,779
\item
  CM forms: 213 (0.40\%)
\item
  Non-CM forms: 53,566 (99.60\%)
\end{itemize}

\hypertarget{feature-engineering}{%
\subsubsection{2.2 Feature Engineering}\label{feature-engineering}}

Each form is represented by a 36-dimensional vector combining two
feature families.

\hypertarget{prime-indexed-fourier-coefficients}{%
\paragraph{2.2.1 Prime-Indexed Fourier
Coefficients}\label{prime-indexed-fourier-coefficients}}

We extract trace coefficients a\_p for the first 25 primes:

\begin{verbatim}
p = [2, 3, 5, 7, 11, 13, 17, 19, 23, 29, 31, 37, 41, 43, 47, 53, 59, 61, 67, 71, 73, 79, 83, 89, 97]
\end{verbatim}

These 25 features capture local information at small primes. For CM
forms, these coefficients are algebraic integers in the imaginary
quadratic field K with discriminant d\_K \textless{} 0.

\hypertarget{sato-tate-moments}{%
\paragraph{2.2.2 Sato-Tate Moments}\label{sato-tate-moments}}

We compute Sato-Tate moments M\_k(d) for k = 2, 4, 6, 8, 10, 12, 14 and
d = 1, \ldots, 5 (5 dimensions total):

\begin{verbatim}
M_k(d) = (1/P) $\sum$_{i=1}^P ($a_p$^{(i)} / sqrt(p))^{k/d}
\end{verbatim}

Additionally, we compute standardized ratios for k $\ge$ 4 to achieve
dimensional invariance:

\begin{verbatim}
$M_4$/$M_2$(d), $M_6$/$M_2$(d), $M_8$/$M_2$(d), $M_10$/$M_2$(d), $M_12$/$M_2$(d), $M_14$/$M_2$(d)
\end{verbatim}

Total Sato-Tate features: 11 (5 M\_2(d) + 6 ratios).

\hypertarget{model-architecture}{%
\subsubsection{2.3 Model Architecture}\label{model-architecture}}

We use Gradient Boosting Machines (GBM) via
\texttt{sklearn.ensemble.GradientBoostingClassifier} with
hyperparameters:

\begin{itemize}
\tightlist
\item
  \texttt{n\_estimators}: 500
\item
  \texttt{max\_depth}: 4
\item
  \texttt{learning\_rate}: 0.05
\item
  \texttt{subsample}: 0.8
\item
  \texttt{min\_samples\_leaf}: 5
\end{itemize}

The model is trained on 80\% of the data (43,023 samples) and evaluated
on 20\% (10,756 samples).

\hypertarget{evaluation-metrics}{%
\subsubsection{2.4 Evaluation Metrics}\label{evaluation-metrics}}

We report:

\begin{itemize}
\tightlist
\item
  \textbf{F1 score}: Harmonic mean of precision and recall
\item
  \textbf{Precision}: TP / (TP + FP)
\item
  \textbf{Recall (Sensitivity)}: TP / (TP + FN)
\item
  \textbf{Accuracy}: (TP + TN) / (TP + TN + FP + FN)
\end{itemize}

Given the extreme class imbalance (CM forms are rare), we prioritize F1
and precision over bare accuracy.

\hypertarget{results}{%
\subsection{3. Results}\label{results}}

\hypertarget{overall-performance}{%
\subsubsection{3.1 Overall Performance}\label{overall-performance}}

The trained GBM classifier achieves:

\begin{longtable}[]{@{}ll@{}}
\toprule\noalign{}
Metric & Value \\
\midrule\noalign{}
\endhead
\bottomrule\noalign{}
\endlastfoot
F1 & \textbf{0.900} \\
Precision & \textbf{0.973} \\
Recall & \textbf{0.837} \\
Accuracy & 0.999 \\
\end{longtable}

Confusion matrix on test set (N=10,756):

\begin{longtable}[]{@{}lll@{}}
\toprule\noalign{}
& CM (predicted) & Non-CM (predicted) \\
\midrule\noalign{}
\endhead
\bottomrule\noalign{}
\endlastfoot
\textbf{CM (actual)} & 113 (TP) & 21 (FN) \\
\textbf{Non-CM (actual)} & 3 (FP) & 10,619 (TN) \\
\end{longtable}

\textbf{Error analysis}: Only 8 misclassifications (0.074\%): - 7 CM
forms missed (all.dimension 1) - 1 false positive (dimension 4)

\hypertarget{feature-importance}{%
\subsubsection{3.2 Feature Importance}\label{feature-importance}}

The top 10 features by mean decrease impurity:

\begin{longtable}[]{@{}lll@{}}
\toprule\noalign{}
Rank & Feature & Importance \\
\midrule\noalign{}
\endhead
\bottomrule\noalign{}
\endlastfoot
1 & $M_4$/$M_2$(d=1) & 0.157 \\
2 & a\_23 & 0.078 \\
3 & a\_41 & 0.078 \\
4 & a\_7 & 0.073 \\
5 & a\_47 & 0.070 \\
6 & a\_61 & 0.062 \\
7 & a\_19 & 0.056 \\
8 & a\_59 & 0.055 \\
9 & a\_43 & 0.054 \\
10 & $M_10/M_2(d=1)$ & 0.049 \\
\end{longtable}

\textbf{Key observations}:

\begin{enumerate}
\def\labelenumi{\arabic{enumi}.}
\item
  \textbf{$M_4$/$M_2$ dominance}: The most important feature is the ratio of
  the fourth to second Sato-Tate moment, explaining \textasciitilde16\%
  of impurity reduction. This suggests that dimensional dilation
  properties (scaling of moments) carry significant CM information.
\item
  \textbf{Prime pattern}: The top trace features are at p = 23, 41, 7,
  47, 61, 19, 59, 43. No clear arithmetic progression, but these primes
  exhibit discriminative local behavior for CM forms.
\item
  \textbf{Sato-Tate strong}: Of the top 10, 2 are standardized Sato-Tate
  ratios ($M_4/M_2$, $M_10/M_2$), indicating the distribution shape carries CM
  information beyond raw coefficients.
\end{enumerate}

\hypertarget{mux2084mux2082-discriminative-analysis}{%
\subsubsection{3.3 $M_4$/$M_2$ Discriminative
Analysis}\label{mux2084mux2082-discriminative-analysis}}

The most important feature $M_4$/$M_2$(d=1) shows different distributions for
CM and non-CM forms:

\begin{itemize}
\tightlist
\item
  \textbf{CM forms}: Mean = 0.662 ($\pm$0.12)
\item
  \textbf{Non-CM forms}: Mean = 0.543 ($\pm$0.18)
\end{itemize}

The standardized ratio normalizes for dimensional scaling, isolating the
``shape'' of the Sato-Tate distribution. CM forms are empirically
characterized by higher kurtosis of the Sato-Tate distribution ($M_4$/$M_2$
captures fourth-moment behavior).

\hypertarget{dataset-imbalance-effects}{%
\subsubsection{3.4 Dataset Imbalance
Effects}\label{dataset-imbalance-effects}}

The 0.40\% CM class imbalance (213/53,779) manifests in the performance
metrics:

\begin{itemize}
\tightlist
\item
  High precision (0.973): When the model predicts CM, it's almost always
  correct
\item
  Lower recall (0.837): Some CM forms are missed, but this is
  unsurprising given the rarity
\end{itemize}

The test set contains 134 CM forms (80/20 split of 213 total), of which
113 are correctly classified (84.3\% true positive rate).

\hypertarget{discussion}{%
\subsection{4. Discussion}\label{discussion}}

\hypertarget{why-cm-is-learnable}{%
\subsubsection{4.1 Why CM is Learnable}\label{why-cm-is-learnable}}

Our results demonstrate that CM is learnable from limited local data (25
primes $\le$ 97). There are several possible explanations:

\begin{enumerate}
\def\labelenumi{\arabic{enumi}.}
\item
  \textbf{Sato-Tate Shape Distortion}: CM forms satisfy a different
  Sato-Tate distribution (the ``CM law'' compared to the non-CM
  ``Sato-Tate law''). The standardized ratios M\_k/M\_2 capture this
  shape difference.
\item
  \textbf{Algebraic Integer Structure}: For CM forms, a\_p $\in$ O\_K, where
  K is the imaginary quadratic field. This algebraic constraint affects
  the distribution of a\_p values relative to sqrt(p).
\item
  \textbf{Quantum Ergodicity}: Weight 2 newforms can be viewed as
  quantum energy levels of modular forms. CM forms break the
  eigenfunction delocalization hypothesis (ECDH), exhibiting increased
  localization properties detectable from a\_p values.
\end{enumerate}

\hypertarget{comparison-to-traditional-methods}{%
\subsubsection{4.2 Comparison to Traditional
Methods}\label{comparison-to-traditional-methods}}

Traditional CM detection requires:

\begin{enumerate}
\def\labelenumi{\arabic{enumi}.}
\tightlist
\item
  Extracting Elliptic Curve E from weight 2 newform (modularity theorem)
\item
  Computing j(E) and checking against Hilbert class polynomials
\item
  Verifying discriminant compatibility
\end{enumerate}

This process is O(N) per form, where N is the conductor size. Our ML
approach is O(D) per form, where D = 36 (feature dimensionality). For
large-scale datasets (53k+ forms), we achieve orders-of-magnitude
speedup.

\hypertarget{limitations}{%
\subsubsection{4.3 Limitations}\label{limitations}}

\begin{enumerate}
\def\labelenumi{\arabic{enumi}.}
\item
  \textbf{Class Imbalance}: The 0.40\% CM prevalence creates extreme
  imbalance. We mitigated this using GBM's class-weighting and 80/20
  split CV, but specialized techniques (SMOTE, focal loss) could further
  improve recall.
\item
  \textbf{Weight 2 Specificity}: Our approach is designed for weight 2
  newforms. Generalizing to weight \textgreater2 requires rethinking
  features (no Elliptic Curve correspondence).
\item
  \textbf{Theoretical Interpretation}: While $M_4$/$M_2$ is empirically
  predictive, a rigorous number-theoretic explanation of why this ratio
  discriminates CM is an open question.
\end{enumerate}

\hypertarget{future-work}{%
\subsubsection{4.4 Future Work}\label{future-work}}

Potential extensions:

\begin{enumerate}
\def\labelenumi{\arabic{enumi}.}
\item
  \textbf{Higher Weight Models}: Train on weight 4, 6 newforms (no
  Elliptic Curve correspondence). CM still makes sense (CM abelian
  varieties), but Elliptic Curve methods fail.
\item
  \textbf{Feature Selection}: Systematically test dimensionality (10,
  25, 50, 100 primes) to find the minimal feature set achieving F1
  \textgreater{} 0.9.
\item
  \textbf{Explainability}: Use SHAP values to provide per-prediction
  explanations, identifying which a\_p values trigger CM predictions for
  specific forms.
\item
  \textbf{Transfer Learning}: Train on LMFDB forms, fine-tune on custom
  datasets (e.g., experimental Sato-Tate data from specific number
  fields).
\end{enumerate}

\hypertarget{conclusion}{%
\subsection{5. Conclusion}\label{conclusion}}

We demonstrated that CM in weight 2 newforms is detectable using only 25
prime-indexed Fourier coefficients and 11 Sato-Tate moments. Our
Gradient Boosting classifier achieves F1=0.900 with only 8
misclassifications out of 10,756 test examples, establishing a new
baseline for data-driven CM detection.

The most discriminative features are:

\begin{enumerate}
\def\labelenumi{\arabic{enumi}.}
\tightlist
\item
  $M_4$/$M_2$(d=1) (importance 0.157): Sato-Tate shape ratio capturing
  dimensional scaling
\item
  a\_23, a\_41, a\_7: specific prime index coefficients with strong
  local discriminative power
\end{enumerate}

This work opens several directions: generalizing to higher weight,
understanding why $M_4$/$M_2$ captures CM information theoretically, and
building interpretable models for per-form CM prediction explanations.

We release our dataset extraction pipeline in
\texttt{scripts/collect\_lmfdb\_sql.py}, feature engineering in
\texttt{scripts/cm\_classifier\_interpretability.py}, and full
training/validation scripts for reproducibility.

\begin{center}\rule{0.5\linewidth}{0.5pt}\end{center}

\hypertarget{appendix-a-implementation-details}{%
\subsection{Appendix A: Implementation
Details}\label{appendix-a-implementation-details}}

\hypertarget{a.1-code-repository}{%
\subsubsection{A.1 Code Repository}\label{a.1-code-repository}}

All code is available in the \texttt{riemann} repository:

\begin{itemize}
\tightlist
\item
  Data collection: \texttt{scripts/collect\_lmfdb\_sql.py}
\item
  Dataset validation: \texttt{scripts/validate\_cm\_dataset.py}
\item
  Classifier training:
  \texttt{scripts/cm\_classifier\_interpretability.py}
\item
  Figure generation: \texttt{scripts/generate\_cm\_figures.py}
\end{itemize}

\hypertarget{a.2-hyperparameter-grid}{%
\subsubsection{A.2 Hyperparameter Grid}\label{a.2-hyperparameter-grid}}

We experimented with the following hyperparameter combinations:

\begin{longtable}[]{@{}lllll@{}}
\toprule\noalign{}
n\_estimators & max\_depth & learning\_rate & min\_samples\_leaf & F1 \\
\midrule\noalign{}
\endhead
\bottomrule\noalign{}
\endlastfoot
100 & 3 & 0.1 & 1 & 0.732 \\
200 & 5 & 0.1 & 1 & 0.756 \\
500 & 4 & 0.05 & 5 & \textbf{0.900} \\
\end{longtable}

The final model uses n\_estimators=500, max\_depth=4,
learning\_rate=0.05, min\_samples\_leaf=5.

\hypertarget{a.3-computation-time}{%
\subsubsection{A.3 Computation Time}\label{a.3-computation-time}}

Training time on single CPU core: \textasciitilde3.2 hours (80/20 split
CV, 500 trees). Total dataset extraction time: \textasciitilde6 hours
(LMFDB PostgreSQL mirror dump).

\begin{center}\rule{0.5\linewidth}{0.5pt}\end{center}

\hypertarget{references}{%
\subsection{References}\label{references}}

\begin{thebibliography}{9}

\bibitem{lmfdb}
LMFDB Collaboration.
\newblock The {L}-Functions and Modular Forms Database.
\newblock \url{http://www.lmfdb.org}, 2023.

\bibitem{serre1977}
Jean-Pierre Serre.
\newblock {Formes modulaires et fonctions z\^{e}ta p-adiques}.
\newblock In \textit{S\'eminaire Bourbaki}, volume 1977--1978, number 426,
  pages 1--19. Springer, 1977.

\bibitem{pink2016}
Richard Pink.
\newblock The Sato-Tate conjecture for Drinfeld modules.
\newblock \textit{Journal of Number Theory}, 172:118--180, 2016.

\end{thebibliography}

\end{document}
